\section{GPU 硬件基础：SM、warp、显存层级、tensor core 与推理瓶颈}

\subsection{为什么第三册必须讲 GPU 硬件}

只讲 GPU API 和 stream 还不够。推理引擎的 GPU 后端最终跑在硬件上：线程如何成组执行，访存如何合并，shared memory 为什么快，tensor core 适合什么形态，显存带宽和 PCIe 传输怎样限制 decode，这些都会决定 AI 编译器 lowering 的选择。如果不了解这些硬件事实，GPU 后端很容易停在“能 launch kernel”的层面。

可以先把 GPU 和第二册里的 CPU 对比。CPU 少量强核心，靠大 cache、乱序执行、分支预测和 SIMD 处理复杂控制流；GPU 大量较简单执行单元，靠海量线程隐藏内存延迟，适合规则、密集、数据并行的工作。两者都要关心内存层级和并行，但并行粒度和瓶颈不同。

\begin{keyidea}
GPU 快不快，取决于工作是否能匹配它的执行模型：大量线程、规则访存、高算术强度、少同步、少 host-device 往返。prefill 往往更匹配，decode 单 token 往往更难匹配。
\end{keyidea}

\subsection{SM/CU：GPU 的基本执行岛}

不同厂商术语不同，NVIDIA 常说 SM，AMD 常说 CU。本册用 SM/CU 表示 GPU 上一组局部执行资源：它包含调度器、寄存器文件、若干向量/标量执行单元、shared memory 或 local data share、load/store 单元，以及可能的 tensor core 或 matrix core。

一个 kernel launch 后，工作会被切成许多 thread block 或 workgroup。每个 block 被调度到某个 SM/CU 上执行。一个 block 内的线程共享一块 shared memory，并能做 block 内同步；不同 block 通常不能直接同步，只能通过全局内存和 kernel 边界间接协作。

\begin{lstlisting}[numbers=none,caption={GPU 执行层级的简化账本}]
kernel launch
  grid: many blocks
    block/workgroup
      warps/wavefronts
        lanes/threads
          registers
      shared memory
  global memory
\end{lstlisting}

这解释了为什么 GPU lowering 要选择 block 形状。一个 attention kernel 可以让一个 block 处理一个 head 的一段 position，也可以让一个 block 处理一组 token 和 head。选择不同，shared memory 用量、寄存器压力、访存合并和 occupancy 都会变。

\subsection{warp/wavefront：线程不是完全独立执行}

GPU 上的线程通常以 warp 或 wavefront 为单位执行。可以把它理解成“一组 lane 同时执行同一条指令，但每个 lane 处理不同数据”。这和 CPU SIMD 有相似性：CPU SIMD 是一条指令处理多个 lane；GPU warp 是多个线程 lockstep 执行同一指令。区别是 GPU 编程模型把 lane 暴露成线程。

如果 warp 内线程走不同分支，会发生 divergence：一部分 lane 先执行分支 A，另一部分 lane 再执行分支 B，吞吐下降。若相邻 lane 访问连续地址，硬件可以合并内存请求；若访问分散，global memory 事务变多。

\begin{center}
\begin{tabular}{p{0.22\linewidth}p{0.30\linewidth}p{0.32\linewidth}}
\toprule
现象 & 类比 & 推理后端影响 \\
\midrule
warp lockstep & SIMD lane 同步执行 & 分支越规则越好 \\
divergence & lane mask 分裂 & sampler/稀疏路径不适合热 kernel \\
coalescing & 连续向量 load & KV cache 和权重 layout 要顺序 \\
occupancy & 同时驻留的 warp 数 & 寄存器/shared memory 过多会降并发 \\
\bottomrule
\end{tabular}
\end{center}

第二册讲 SIMD 时强调“相邻 lane 和规则控制流”。GPU 这里是同一原则的放大版。prefill 的大矩阵乘通常规则；decode 中某些小 batch、变长 context、条件路径和 host 同步更容易破坏规则性。

\subsection{coalescing：连续地址决定一次访存能带回多少有效数据}

GPU global memory 带宽很高，但前提是 warp/wavefront 里的 lane 访问足够连续。若 32 个 lane 分别访问相邻的 fp16 元素，硬件可以用较少内存事务带回一整段数据；若 32 个 lane 访问相隔很远的位置，就会产生更多事务，实际带宽下降。

\begin{lstlisting}[numbers=none,caption={coalesced 与 scattered 访问的差异}]
coalesced:
  lane 0 reads base + 0
  lane 1 reads base + 1
  lane 2 reads base + 2
  ...

scattered:
  lane 0 reads page_table[a]
  lane 1 reads page_table[b]
  lane 2 reads page_table[c]
  ...
\end{lstlisting}

KV cache layout 直接影响 coalescing。若一个 warp 负责同一个 head 的连续 \code{head_dim}，那么 \code{head_dim} 连续布局有利；若一个 warp 负责多个 position 的同一维度，那么 position 连续更有利。paged KV 会引入 page table，长上下文和多 session 调度更灵活，但 kernel 必须小心把 page 内连续读和跨页边界分开处理。

\begin{warningbox}
GPU 上“逻辑上连续”不等于“lane 访问连续”。descriptor 只写 shape 不够，lowering 必须知道 warp mapping 和物理 stride。
\end{warningbox}

\subsection{显存层级：global、shared、register 和 cache}

GPU 的 global memory 容量大、带宽高，但延迟也高。shared memory 在 SM/CU 内部，容量小但速度快，适合 tile 复用。register 是每个线程私有，最快但数量有限。L2 cache 连接多个 SM/CU，L1 或 texture cache 细节随架构不同。

\begin{systemdiagram}{GPU 推理中的内存层级}
\begin{pdfdiagram}[x=1cm,y=1cm]
  \node[book accent node,text width=2.2cm] (host) at (0,1.9) {host memory\\CPU};
  \node[book teal node,text width=2.2cm] (global) at (3.0,1.9) {global memory\\HBM/GDDR};
  \node[book purple node,text width=2.2cm] (l2) at (6.0,1.9) {L2/cache\\跨 SM};
  \node[book amber node,text width=2.2cm] (shared) at (9.0,2.75) {shared\\memory};
  \node[book navy node,text width=2.2cm] (reg) at (9.0,1.05) {registers\\per thread};

  \draw[book strong arrow] (host) -- node[book label,above]{PCIe/NVLink} (global);
  \draw[book strong arrow] (global) -- (l2);
  \draw[book strong arrow] (l2) -- (shared);
  \draw[book strong arrow] (l2) -- (reg);
\end{pdfdiagram}
\diagramnote{本图要证明的命题是：GPU kernel 性能不只取决于算力，还取决于数据在哪一层内存里移动。}
\end{systemdiagram}

对推理来说，权重和 KV cache 通常在 global memory；tile 化 GEMM 会把输入/权重片段搬到 shared memory 或寄存器；accumulator 放在寄存器；logits 可能需要 staging 到 host。若每个 decode step 都把完整 logits 从 global memory 拷回 host，PCIe 或同步开销可能比某个小 kernel 更显眼。

\subsection{shared memory：它解决的是复用，不是容量}

shared memory 适合保存一个 block 内多个线程都会反复使用的小 tile。矩阵乘中，一个 input tile 和 weight tile 被多个 lane 复用，搬到 shared memory 后可以减少 global memory 读流。attention 中，也可以把 K/V 的一段 block 或 softmax 中间量放到 shared memory，减少重复读写。

但 shared memory 容量小，而且占用越多，每个 SM/CU 能同时驻留的 block 可能越少。用它之前要问三个问题：

\begin{itemize}
  \item 这个 tile 会被多少次复用，是否值得从 global memory 搬进来。
  \item tile 大小是否会让 occupancy 明显下降。
  \item shared memory 的访问是否会发生 bank conflict。
\end{itemize}

例如 decode attention 读取历史 K/V 时，如果每个 K/V 元素只被一个 query head 用一次，把整段历史搬进 shared memory 可能不划算；若同一 K/V block 会被多个 query head 或多个 batch row 复用，就有价值。GQA 的共享 K/V head 也会改变这个判断：多个 query head 读同一个 K/V head，理论上有复用机会，但具体能否利用取决于 kernel mapping。

\subsection{tensor core：不是所有矩阵都天然吃满}

现代 GPU 常有 tensor core 或 matrix core，专门加速小矩阵块乘加。它们对 dtype、矩阵形状、对齐和 layout 有要求，常见于 fp16/bf16/tf32/int8/fp8 等路径。prefill 的大 GEMM 更容易使用 tensor core；decode 单 token 的 GEMV 不一定能有效使用。

\begin{center}
\begin{tabular}{p{0.22\linewidth}p{0.30\linewidth}p{0.32\linewidth}}
\toprule
场景 & tensor core 适配性 & 关键条件 \\
\midrule
prefill QKV/MLP GEMM & 通常较好 & batch/token 足够大，layout 对齐 \\
decode GEMV & 通常较难 & M 维太小，launch 和访存占比高 \\
int4 权重 & 取决于格式支持 & packed layout 和 scale 处理 \\
attention & 取决于 kernel 设计 & softmax、mask、KV 读流复杂 \\
\bottomrule
\end{tabular}
\end{center}

这解释了一个常见现象：GPU prefill 很快，但 decode tokens/s 未必按峰值算力提升。decode 阶段每次只有一个 token，矩阵形态小，KV cache 读历史，logits 还要进 sampler。算力不是唯一瓶颈。

\subsection{用 roofline 判断 GPU 推理瓶颈}

roofline 的核心是算术强度：

\[
AI=\frac{FLOPs}{Bytes}
\]

若 \code{AI} 很低，即使 GPU 峰值 FLOPS 很高，也会被显存带宽限制；若 \code{AI} 很高，才更可能接近计算单元上限。prefill GEMM 的权重 tile 被多个 token 行复用，算术强度高；decode GEMV 的 batch 通常很小，权重复用弱，算术强度低。KV attention 还要随上下文读取历史 K/V，常常更像带宽问题。

\begin{center}
\begin{tabular}{p{0.18\linewidth}p{0.34\linewidth}p{0.32\linewidth}}
\toprule
路径 & 算术强度直觉 & 常见瓶颈 \\
\midrule
prefill GEMM & 高，tile 复用好 & tensor core、workspace、activation \\
decode GEMV & 低到中，权重复用弱 & 显存带宽、launch、小矩阵效率 \\
decode attention & 随 context 读 KV & KV 布局、带宽、softmax 归约 \\
logits staging & 计算少，传输和同步明显 & host-device copy、sampler boundary \\
\bottomrule
\end{tabular}
\end{center}

一个简单反证方法是先算字节下限。若某个 decode step 至少要读 4GB 权重和 KV，而有效显存带宽是 800GB/s，那么单步光读数据就至少 5ms，还没算 kernel launch、softmax、MLP、staging 和 sampler。若 benchmark 报告低于这个数量级，就要检查是否统计了完整 step、是否用了 batch、是否跳过了某些层、是否缓存命中口径不同。

\begin{keyidea}
GPU 性能分析不能只问“有没有 tensor core”。先算 FLOPs、bytes、launch、copy、sync，再判断瓶颈更像计算、带宽、调度还是 host-device 边界。
\end{keyidea}

\subsection{occupancy：并发多不等于一定快}

occupancy 指一个 SM/CU 上同时驻留的 warp/block 相对硬件上限的比例。高 occupancy 可以帮助隐藏 global memory 延迟，但不是越高越好。寄存器用太多、shared memory 用太多、block 太大，都会降低 occupancy；反过来，为了提高 occupancy 把 tile 做得太小，也可能降低数据复用。

推理 kernel 要平衡：

\begin{itemize}
  \item 寄存器：accumulator 多，复用好，但 occupancy 可能降。
  \item shared memory：tile 大，global memory 读少，但可驻留 block 少。
  \item block 大小：并行度高，但同步和资源压力上升。
  \item 算术强度：每读一个字节做多少计算，决定更像算力瓶颈还是带宽瓶颈。
\end{itemize}

因此 GPU 后端的 profiler 不应只报告 occupancy。一个 kernel occupancy 很高但访存不合并，仍然慢；一个 tensor core kernel occupancy 不高但计算吞吐高，也可能很快。指标要和具体瓶颈一起看。

\subsection{warp 映射和 tensor core 账本}

GPU kernel 设计必须写清 warp mapping：lane 维度对应 token、head、head\_dim、output channel 还是 K tile。不同映射会改变 coalescing、shared memory 复用、reduction 方式和 divergence。只写 \code{blockDim=256} 没有解释力。

\begin{center}
\begin{tabular}{p{0.24\linewidth}p{0.34\linewidth}p{0.28\linewidth}}
\toprule
warp mapping & 优点 & 风险 \\
\midrule
lane 覆盖连续 \code{head\_dim} & K/V 读连续，coalescing 好 & head\_dim 小时并行度有限 \\
lane 覆盖 position & 长上下文并行度高 & softmax reduction 和 mask 更复杂 \\
warp 覆盖 output channels & GEMV 权重读规则 & batch=1 时复用弱 \\
warp 覆盖 batch rows & 权重复用好 & 小 batch 或变长 session 难凑齐 \\
\bottomrule
\end{tabular}
\end{center}

\topic{shared memory bank conflict}
shared memory 虽然快，但不是“无成本数组”。它由多个 bank 组成，若一个 warp 的多个 lane 同时访问同一 bank 的不同地址，访问会被串行化或降效。Tile layout 要同时满足 global memory coalescing 和 shared memory bank 访问。一个从 global 读连续的布局，搬到 shared 后若按错误 stride 读取，仍可能慢。

\topic{tensor core 的形状合同}
Tensor core 通常执行固定小矩阵块的乘加，例如某种 \code{m x n x k} tile。后端要把大 GEMM 切成这些 tile，并满足 dtype、对齐、layout 和 accumulator 要求。若 decode 是 \code{M=1} 的 GEMV，很多 tensor core tile 的 M 维利用率很低；把多个 session 合批可以提高 M，但会引入 queue wait 和 KV binding 复杂度。

\begin{lstlisting}[numbers=none,caption={tensor core 使用前的账本}]
matrix shape:
  M = batch_or_tokens
  N = output_channels
  K = hidden

requirements:
  dtype supported by matrix core
  M/N/K aligned to tile or tail handled
  input/weight layout matches fragment loader
  accumulator dtype and scale path defined
  batch wait cost included in SLO
\end{lstlisting}

\topic{register pressure 和 occupancy 的交换}
更大的 tile 会增加寄存器 accumulator，可能提升复用，也可能让每个 SM 驻留 warp 变少。更高 occupancy 能隐藏内存延迟，但如果每个 warp 都在读不合并的 global memory，occupancy 只是让更多慢请求同时排队。GPU 报告要同时写 registers per thread、shared memory per block、achieved occupancy、memory throughput 和 tensor core utilization。

\subsection{GPU profiler 指标怎样读}

不同工具名字不同，但读法相近。不要孤立看一个百分比，而要把指标放回 kernel 目标。

\begin{center}
\begin{tabular}{p{0.24\linewidth}p{0.32\linewidth}p{0.28\linewidth}}
\toprule
指标 & 可能说明什么 & 需要结合什么看 \\
\midrule
achieved occupancy & 驻留 warp 是否足够 & 寄存器、shared memory、block size \\
memory throughput & 是否接近带宽上限 & coalescing、cache hit、字节估算 \\
tensor core utilization & 矩阵单元是否忙 & dtype、layout、tile、GEMM 形态 \\
warp stall & 等内存、同步或依赖 & global load、barrier、寄存器依赖 \\
launch count & 小 kernel 数量 & decode step latency、fusion 机会 \\
copy time & host-device 传输 & staging 大小、sampler 位置 \\
\bottomrule
\end{tabular}
\end{center}

例如 memory throughput 很高而 tensor core utilization 低，不一定是坏事；decode attention 可能本来就是读 KV 的带宽路径。反过来，occupancy 很高但 throughput 很低，可能是访问不合并、分支发散或 page table 间接访问导致。报告必须把这些指标和 executable plan 中的 kernel desc 对上。

\subsection{prefill 和 decode 的 GPU 瓶颈不同}

把前面硬件事实放回推理链路，可以得到一个清晰判断：

\begin{center}
\begin{tabular}{p{0.18\linewidth}p{0.34\linewidth}p{0.32\linewidth}}
\toprule
阶段 & GPU 优势 & 主要风险 \\
\midrule
prefill & 大 GEMM、并行 token、tensor core & activation 峰值、attention workspace、显存占用 \\
decode & 权重常驻、KV cache 常驻、可融合小算子 & launch、同步、GEMV 小形态、KV 读流、logits staging \\
采样 & 可做 top-k 或局部 sampler & 完整 logits 回传和 CPU 同步 \\
量化 & 降低显存和带宽 & 格式不支持、scale metadata、dequant 开销 \\
\bottomrule
\end{tabular}
\end{center}

这张表是 GPU lowering 的基础。prefill plan 应优先选择高吞吐 GEMM/attention；decode plan 应优先减少 launch 和 host-device 同步，优化 KV cache layout，并按 context length 分档。若把 prefill 的大矩阵思路直接套到 decode，通常会失望。

\subsection{GPU 硬件事实怎样进入 AI 编译器}

GPU 硬件不应该散落成后端里的魔法数字。它应该进入 capability 和 kernel descriptor：

\begin{lstlisting}[numbers=none,caption={GPU capability 和 kernel descriptor 字段}]
GpuCapability:
  backend_name
  max_threads_per_block
  warp_or_wavefront_size
  shared_memory_per_block
  register_file_limits
  supported_matrix_dtypes
  supported_quant_formats
  host_device_link

GpuKernelDesc:
  operator
  stage: prefill | decode
  block_shape
  warp_mapping
  shared_memory_bytes
  registers_estimate
  supported_layouts
  sync_points
\end{lstlisting}

编译器 lowering 使用这些字段选择 plan；verifier 用它们拒绝不支持的 shape 或 dtype；profiler 用它们解释瓶颈；benchmark 报告用它们保证可复现。这样 GPU 硬件知识就进入了工程合同，而不是停留在“这个 kernel 在我机器上更快”的经验。

\subsection{coalescing 的地址公式：warp 访问不是自然合并}

GPU global memory 喜欢一个 warp 的 lane 访问连续、对齐、少量内存事务能覆盖的地址。若 lane \(i\) 访问：

\[
addr_i = base + (offset + i) \times sizeof(T)
\]

通常更容易合并。若 lane \(i\) 访问：

\[
addr_i = base + index_i \times sizeof(T)
\]

且 \(index_i\) 随机，硬件仍要发出多个内存事务。coalescing 不是“用了 GPU 就自动发生”，它来自 layout、warp mapping 和地址公式。

\begin{lstlisting}[numbers=none,caption={KV cache 两种 warp 地址流}]
good for lane covers head_dim:
  lane i reads kv[layer][pos][kv_head][dim + i]
  addresses are contiguous inside head_dim

bad if lane covers random sessions:
  lane i reads kv[session_i][page_i][slot_i][dim]
  addresses may scatter across pages
\end{lstlisting}

paged KV cache 提高调度灵活性，但会让地址公式多一次 page table 间接。高质量 kernel 通常要把 page 内连续读和跨页边界分开处理，避免每个 lane 都走完全随机路径。

\subsection{shared memory bank conflict 的手算模型}

shared memory 被拆成多个 bank。简化地说，若有 \(B\) 个 bank，地址按 word 映射到：

\[
bank = word\_index \bmod B
\]

若一个 warp 中多个 lane 同时访问同一 bank 的不同地址，就可能冲突。连续访问通常较好；按某个 stride 访问可能让所有 lane 落到少数 bank。

\begin{lstlisting}[numbers=none,caption={bank conflict 的 stride 直觉}]
bank_count = 32

lane i reads shared[i]:
  bank = i mod 32
  lanes spread across banks

lane i reads shared[i * 32]:
  bank = 0 for all lanes
  worst conflict shape
\end{lstlisting}

实际硬件有广播、bank 宽度和访问粒度等细节，但这个模型能解释为什么 shared memory tile layout 需要 padding 或 swizzle。把 global memory 读连续搬进 shared 只是第一步；shared 内部消费也要看 bank。

\subsection{occupancy calculator 的最小公式}

occupancy 受 threads per block、registers per thread、shared memory per block 和硬件上限共同限制。简化账本如下：

\begin{lstlisting}[numbers=none,caption={occupancy 限制项}]
blocks_by_threads = max_threads_per_sm / threads_per_block
blocks_by_registers = registers_per_sm /
  (registers_per_thread * threads_per_block)
blocks_by_shared = shared_memory_per_sm / shared_memory_per_block
resident_blocks = min(blocks_by_threads,
                      blocks_by_registers,
                      blocks_by_shared,
                      hardware_block_limit)
\end{lstlisting}

优化时常见交换是：增大 tile 提高复用，但增加 registers/shared，降低 resident blocks；减小 tile 提高 occupancy，但每个 global load 被复用次数下降。报告里只写 occupancy 百分比不够，必须写是哪一项限制了驻留。

\begin{center}
\begin{tabular}{p{0.24\linewidth}p{0.34\linewidth}p{0.28\linewidth}}
\toprule
限制项 & 表面现象 & 可能改法 \\
\midrule
registers/thread 高 & resident warp 少，可能 spill & 缩小 tile，分阶段累加 \\
shared/block 高 & resident block 少 & 压缩 tile，重用 shared buffer \\
threads/block 不合适 & 并行度或调度粒度差 & 调整 block shape \\
memory coalescing 差 & throughput 低、stall memory & 改 warp mapping/layout \\
\bottomrule
\end{tabular}
\end{center}

\subsection{decode 小 batch 为什么难吃满 GPU}

LLM decode 常见痛点是 batch 小、每步依赖上一步 token、kernel 多、同步多。单 session decode 的 GEMV 形态 \(M=1\)，很难像 prefill GEMM 那样提供大矩阵给 tensor core。即使权重常驻显存，每一步仍要调度多个 kernel，读大量权重和 KV，最后 sampler 可能还要同步 logits。

\begin{lstlisting}[numbers=none,caption={单 session decode 的 GPU 固定成本}]
for each token:
  launch or dispatch many small kernels
  read layer weights
  read KV cache up to context
  run reductions for attention
  produce logits
  synchronize sampler or copy top candidates
  feed next token to next step
\end{lstlisting}

连续合批能提高 \(M\)，让 GEMV 更像小 GEMM，也能摊薄 launch 成本；代价是 queue wait、KV binding、变长 session 和 SLO 复杂度。GPU backend 不能只追求吞吐，它必须和第三章的 scheduler 一起设计。

\subsection{host-device 边界和 sampler 位置}

如果每步把完整 logits 从 GPU 拷回 CPU 再采样，词表大小 \(V\) 很大时会产生明显传输和同步成本。更好的路径可能是在 GPU 上做 logits processor、top-k/top-p 或至少取 top candidates，再只传小结果；但这会把采样语义、RNG 状态和可复现性带进 GPU plan。

\begin{center}
\begin{tabular}{p{0.24\linewidth}p{0.34\linewidth}p{0.28\linewidth}}
\toprule
采样位置 & 优点 & 风险 \\
\midrule
CPU 完整 logits & 易调试，复现简单 & 大拷贝、同步、TPOT 增加 \\
GPU top-k/top-p & 传输少，减少同步 & kernel 复杂，RNG/排序合同 \\
混合 top candidates & 折中，保留 CPU sampler & 需要误差和截断边界 \\
\bottomrule
\end{tabular}
\end{center}

这也是为什么 executable plan 要记录 sampler boundary。否则 benchmark 的 TPOT 可能测的是 GPU kernel，而用户体验卡在 logits staging 和 CPU sampler。

\subsection{GPU 报告的工业级最小字段}

GPU benchmark 报告至少包含：

\begin{lstlisting}[numbers=none,caption={GPU backend report 最小字段}]
model_shape
stage: prefill | decode
batch_size
context_length
kernel_count_per_token
selected_kernels
block_shape
warp_mapping
registers_per_thread
shared_memory_per_block
achieved_occupancy
memory_throughput
tensor_core_utilization
host_device_copy_bytes
sampler_boundary
correctness_oracle
\end{lstlisting}

这些字段让 GPU 优化回到证据链：kernel 是否真的对应目标阶段，shape 是否足够吃满硬件，访存是否合并，tensor core 是否被使用，host-device 是否拖尾，结果是否仍对。

\section{状态机：一个 GPU kernel 怎样消耗资源}

GPU 章节不能只背 thread、block、warp、SM。一次 kernel 从 launch 到完成，也是一台资源状态机：

\begin{center}
\begin{tabular}{p{0.20\linewidth}p{0.34\linewidth}p{0.32\linewidth}}
\toprule
状态 & 资源账本 & 典型瓶颈 \\
\midrule
Launched & grid/block、参数、stream & launch overhead、依赖同步 \\
Resident & block 占用 SM，分配寄存器/shared memory & occupancy 受寄存器或 shared 限制 \\
WarpScheduled & warp 被调度发射指令 & divergence、指令供给 \\
MemoryIssued & global/shared/register 访问发出 & coalescing、bank conflict、带宽 \\
ComputeIssued & ALU/tensor core 指令发出 & tile shape、dtype、数据准备 \\
Writeback & 结果写回 global 或下一 kernel & store 合并、同步、拷贝 \\
Completed & event 可见，后续依赖解除 & host-device 同步拖尾 \\
\bottomrule
\end{tabular}
\end{center}

这张表让 GPU 性能报告不再只写“occupancy 低”。occupancy 低可能是寄存器太多，也可能是 shared memory 太多；occupancy 高也可能因为访存不合并而慢。必须把 resident blocks、warp scheduling、memory path 和 compute path 分开解释。

\subsection{数学模型：occupancy 不是越高越好}

设每个 SM 最大线程数为 \(T_{max}\)，最大 block 数为 \(B_{max}\)，寄存器总数为 \(R_{sm}\)，shared memory 为 \(S_{sm}\)。若每个 block 有 \(T_b\) 个线程，每线程用 \(R_t\) 个寄存器，每 block 用 \(S_b\) shared memory，则 resident block 上限近似为：

\[
B_{res} \le \min\left(B_{max}, \left\lfloor\frac{T_{max}}{T_b}\right\rfloor,
\left\lfloor\frac{R_{sm}}{T_b R_t}\right\rfloor,
\left\lfloor\frac{S_{sm}}{S_b}\right\rfloor\right)
\]

occupancy 是可驻留 warp 与硬件最大 warp 的比值。它的作用是隐藏 latency；但若 kernel 主要受 global memory 带宽限制，或者寄存器减少导致 spill，盲目提高 occupancy 可能更慢。GPU 报告要同时列 occupancy、register spill、memory throughput 和 achieved compute。

\subsection{coalescing 和 bank conflict 的具体反例}

一个 warp 中第 \(i\) 个 lane 访问地址：

\[
addr_i = base + (lane_i \times stride) \times sizeof(T)
\]

当 stride=1 且地址对齐时，访问容易合并；当 stride 很大时，一个 warp 可能触碰许多 cache line 或 memory transaction。shared memory 也类似，若多个 lane 落到同一个 bank，访问会序列化或产生冲突。

\begin{lstlisting}[numbers=none,caption={GPU 访存事故复盘字段}]
symptom:
  kernel occupancy acceptable but memory throughput low

address model:
  lane i reads base + i * stride * sizeof(float)
  stride = hidden_size
  warp touches many cache lines

counter evidence:
  low global load efficiency or fallback timing evidence
  shared bank conflicts if available

fix:
  change layout or warp mapping
  tile into shared memory
  make lane contiguous on innermost dimension
\end{lstlisting}

这个反例对应 AI 里的真实问题：tensor layout 若让 warp lane 跨 head 或跨 token 跳跃访问，GPU 后端会把显存带宽浪费在不合并访问上。layout propagation 不是 compiler 装饰，而是 GPU 性能条件。

\subsection{工业取舍：tensor core、launch 和可复现性}

tensor core 适合形状足够大、dtype/layout 符合要求、tile 能铺满的矩阵。prefill 常能满足；单 session decode 常常是小 batch GEMV，利用率差。把多个 session 合批可以提高 tile 利用，但会引入排队等待、KV gather、变长 mask 和 SLO 取舍。若把 sampler 也搬到 GPU，还要处理 RNG 状态、top-k/top-p 排序稳定性和跨设备复现。

\begin{center}
\begin{tabular}{p{0.24\linewidth}p{0.32\linewidth}p{0.30\linewidth}}
\toprule
优化 & 收益 & 新风险 \\
\midrule
continuous batching & 提高矩阵形状和吞吐 & queue wait、KV 绑定复杂 \\
GPU sampler & 减少 logits 拷贝和同步 & RNG/排序/复现合同复杂 \\
fused kernel & 降低 launch 和读写中间值 & debug/oracle materialization 变难 \\
tensor core lowering & 高吞吐 & dtype/layout/tile 限制，误差边界 \\
\bottomrule
\end{tabular}
\end{center}

\subsection{实验：GPU 报告必须有反证输入}

\begin{exercisebox}
实验一：设计 prefill 与 decode 两组 shape。prefill 使用 \(T=128\) 或更大，decode 使用 \(T=1\)。记录 kernel count、host-device copy bytes、batch size、context length 和 TPOT/TTFT。解释为什么两者不能用同一个 tokens/s 结论。

实验二：构造 stride=1 和 stride=hidden 的两种 lane 映射。即使没有 GPU，也要写出地址公式、cache line transaction 预期和应采集的 profiler 字段。

实验三：比较 CPU sampler、GPU top-k 和混合 top candidates 三种边界。固定 logits、seed 和 sampler state，说明如何验证 token 一致或误差可接受。

实验四：填写一次 GPU kernel 的 occupancy 账本。至少包含 block size、registers/thread、shared memory/block、resident blocks、achieved occupancy 和是否发生 spill。
\end{exercisebox}

\subsection{本节验收线}

读完本节，应该能把 GPU 硬件和推理后端连接起来：

\begin{itemize}
  \item SM/CU 是 GPU 的局部执行资源，block/workgroup 被调度到其上。
  \item warp/wavefront 类似放大的 SIMD，喜欢规则控制流和连续访存。
  \item global memory、shared memory、register、cache 和 host-device 链路共同决定数据移动成本。
  \item tensor core 适合满足 dtype/layout/shape 条件的大矩阵块，不自动解决 decode GEMV。
  \item occupancy 要和寄存器、shared memory、访存合并、算术强度一起看。
  \item prefill 和 decode 的 GPU 瓶颈不同，lowering 和 profiler 必须分开。
  \item GPU capability 应进入 descriptor、verifier、lowering 和报告。
\end{itemize}
